\begin{references}
\item Manson JE, Cook NR, Lee I-M, et al. Vitamin D supplements and prevention of cancer and cardiovascular disease. N Engl J Med. 2019;380:33-44.
\item Manson JE, Cook NR, Lee I-M, et al. Marine n−3 fatty acids and prevention of cardiovascular disease and cancer. N Engl J Med. 2019;380:23-32.
\item Dumoulin C, Cacciari LP, Hay-Smith EJC. Pelvic floor muscle training versus no treatment, or inactive control treatments, for urinary incontinence in women. Cochrane Database Syst Rev. 2018;10(10):CD005654.
\item Dumoulin C, Morin M, Danieli C, et al., for the Urinary Incontinence and Aging Study Group. Group-based vs individual pelvic floor muscle training to treat urinary incontinence in older women: a randomized clinical trial. JAMA Intern Med. 2020;180(10):1284-1293.
\item Piaggio G, Elbourne DR, Altman DG, et al. Reporting of non-inferiority and equivalence randomized trials. An extension of the CONSORT Statement. JAMA. 2006;295:1152-1160.
\item Piaggio G, Elbourne DR, Pocock SJ, et al. Reporting of non-inferiority and equivalence randomized trials. An extension of the CONSORT 2010 statement. JAMA. 2012;308:2594-2604.
\item D’Agostino RB Sr., Massaro JM, Sullivan LM, et al. Non-inferiority trials: design concepts and issues—the encounters of academic consultants in statistics. Statist Med. 2003;22:169-186.
\item Kaul S, Diamond GA. Good enough: a primer on the analysis and interpretation of non-inferiority trials. Ann Intern Med. 2006;145:62-69.
\item Chang M, Chow S, Pong A. Adaptive design in clinical research: issues, opportunities, and recommendations. J Biopharm Stat. 2006;16:299-309.
\item Berry SM, Connor JT, Lewis RJ. The Platform Trial: an efficient strategy for evaluating multiple treatments. JAMA. 2015;313:1619-1620.
\item Beigle JH, Tomashek KM, Dodd LE, et al., for the ACTT-1 Study Group Members. Remdesivir for the treatment of Covid-19—final report. N Engl J Med. 2020;383:1813-1826.
\item Guyatt G, Sackett D, Taylor DW, et al. Determining optimal therapy— randomized trials in individual patients. N Engl J Med. 1986:314:889-892.
\item Walsh M, Hilton J, Masouredis C, et al. Smokeless tobacco cessation intervention for college athletes: results after 1 year. Am J Public Health. 1999;89:228-234.
\item Donner A, Birkett N, Buck C, et al. Randomization by cluster: sample size requirements and analysis. Am J Epidemiol. 1981;114:906-914.
\item Cook AJ, Delong E, Murray DM, et al. Statistical lessons learned for designing cluster randomized pragmatic clinical trials from the NIH health care systems collaboratory biostatistics and design core. Clin Trials. 2016;13:504-512.
\item Hemming K, Haines TP, Chilton PJ, et al. The stepped wedge cluster randomised trial: rationale, design, analysis, and reporting. BMJ. 2015;350:h391.
\item Ellenberg SS. The stepped-wedge clinical trial. Evaluation by rolling deployment. JAMA. 2018;319:607-608.
\item Gruber JS, Reygadas F, Arnold BF, Ray I, Nelson K, Colford JM. A stepped wedge, cluster-randomized trial of a household UV-disinfection and safe storage drinking water intervention in rural Baja California Sur, Mexico. Am J Trop Med Hyg. 2013;89(2):238-245.
\item Chalmers T, Celano P, Sacks H, et al. Bias in treatment assignment in controlled clinical trials. N Engl J Med. 1983;309:1358-1361.
\item Sedgwick P. Before and after study designs. BMJ. 2014;349:g5074.
\item Bernal JL, Cummings S, Gasparrini A. Interrupted time series regression for the evaluation of public health interventions: a tutorial. Int J Epidemiol. 2017;46:348-355.
\end{references}
